Dans un modèle à une couche plus réaliste, l'atmosphère absorbe seulement une fraction $\varepsilon = \num{\v{eps}}$ du rayonnement infrarouge émis par le sol. On peut montrer que la température du sol est alors
\[ T_{\text{sol}} = T_e\left(\frac{2}{2 - \varepsilon}\right)^{1/4} \]
où $T_e = \qty{\v{Te}}{K}$ est la température d'équilibre de la planète sans atmosphère.

\begin{enumerate}[label=\alph*)]
    \item Vérifiez que cette formule redonne les cas limites $\varepsilon = 0$ (pas d'effet de serre) et $\varepsilon = 1$ (atmosphère totalement opaque à l'infrarouge).
    \item Déterminez la température du sol pour $\varepsilon = \num{\v{eps}}$.
    \item Une hausse de la concentration des gaz à effet de serre fait passer $\varepsilon$ à \num{\v{eps2}}. De combien la température du sol augmente-t-elle ?
\end{enumerate}
